% SNAS 2026 standalone abstract. Text is kept identical to the abstract in
% SNAS_2026_FaithBench_short_paper.tex; update both together.
% NOTE: the "_blind" filename is historical. This now carries the camera-ready
% abstract. Author names still live only in the paper, not here.
\documentclass[12pt,letterpaper,twoside]{article}

\usepackage[margin=1in]{geometry}
\usepackage{fontspec}
\setmainfont{Times New Roman}
\usepackage{setspace}
\doublespacing
\usepackage{fancyhdr}
\usepackage[hidelinks]{hyperref}
\hypersetup{
  pdftitle={ConstructionSafety-FaithBench Abstract},
  pdfauthor={},
  pdfsubject={Double-blind abstract submission candidate},
  pdfkeywords={trustworthy AI, explainable AI, vision-language models, construction safety}
}

\setlength{\parindent}{0.5in}
\setlength{\emergencystretch}{3em}
\pagestyle{fancy}
\fancyhf{}
\fancyhead[LE,LO]{\textit{ConstructionSafety-FaithBench}}
\fancyhead[RE,RO]{\thepage}
\renewcommand{\headrulewidth}{0.4pt}
\setlength{\headheight}{15pt}

\begin{document}

\begin{center}
\textbf{ConstructionSafety-FaithBench: Evaluating Visual-Evidence Sensitivity in Construction-Safety Vision-Language Systems}
\end{center}

\begin{center}
\textbf{Abstract}
\end{center}

Construction safety depends on small, rule-specific visual details: a hard hat on a particular worker, a barrier at a particular edge. Vision-language models (VLMs) are increasingly proposed to check such details automatically, but they can give the right safety answer for the wrong visual reason, judging a scene unsafe because construction sites look hazardous while ignoring the cue the rule concerns. Existing evaluations report answer accuracy and are silent on this failure. We introduce ConstructionSafety-FaithBench, an audit benchmark that separates three properties usually merged into one score: whether the answer is correct, whether the cited visual evidence is rule-relevant, and whether the answer changes when that evidence is removed. Evidence dependence is probed by occluding the rule-relevant region and comparing the effect against occlusions of identical size placed elsewhere in the image, so any difference reflects location rather than the amount removed. The benchmark is model-agnostic; here we validate it on one system, an open-vocabulary grounding pipeline built on Florence-2, across hard-hat, fall-protection, guardrail, and struck-by rules. For that system the three properties come apart. The pipeline reaches 18.3\% accuracy on 120 adjudicated hard cases. It returns an evidence region on 98.8\% of images, yet those regions overlap the rule-relevant one poorly, mean intersection-over-union 0.058. It nonetheless flips its answer 39.2\% of the time under targeted occlusion, against 9.2\% under size-matched random occlusion, a paired difference of 30.0 points. It never answers ``uncertain,'' even where a single frame cannot support a decision. Coming from one system on one imagery source, these findings characterize that pipeline rather than vision-language models in general. They nonetheless indicate that answer accuracy is insufficient evidence of trustworthiness for safety-critical screening, and that a system's willingness to abstain deserves scrutiny alongside its accuracy. Benchmark, code, and frozen artifacts: \url{https://github.com/kamrul28890/snas-2026-construction-safety-faithbench}.

\end{document}
